\documentclass[10pt]{article}

% -- Packages --
\usepackage[margin=1in]{geometry}
\usepackage{graphicx}
\usepackage{listings}
\usepackage{amsmath,amssymb}
\usepackage{url}

\title{CS2650 Midway Check-in Design Document}
\author{Nico Fidalgo \\
Project: LSM-Tree Key-Value Store}
\date{March 26, 2025}

\begin{document}
\maketitle

\section{Introduction}
Briefly describe your project goals and the context of building
an LSM-tree-based key-value store. State your primary focus:
an in-memory buffer, on-disk levels, merges, and dictionary 
operations supporting \texttt{put}, \texttt{get}, \texttt{range}, \texttt{delete}, \texttt{load}.
Explain you are not yet focusing on concurrency.

\section{Data Layout and Organization}
\subsection{In-memory Buffer (Level-0)}
Discuss the data structure (e.g., skip list). 
Explain the flush threshold and how data is sorted and flushed.

\subsection{On-Disk Levels}
Detail how runs (SSTables) are stored, file format, 
level capacities, and how you keep track of key ranges. 
Note any fence pointer approach, Bloom filters usage or plan.

\subsection{Merging / Compaction}
Describe your chosen merge policy (tiering, leveling, or both),
how merges are triggered, and how tombstones are processed. 
Mention I/O complexities briefly.

\section{Dictionary Operations}
\subsection{Put/Update}
Pseudo-code or a short figure to show how inserts flow 
to memtable and eventually to disk.

\subsection{Get}
Summarize how a lookup goes from top to bottom (memtable, Bloom filters, etc.).
Brief pseudo-code for clarity.

\subsection{Range}
Explain the method for scanning multiple runs and merging results.

\subsection{Delete}
Note that deletes place a tombstone in the memtable and merges eventually discard them.

\section{Preliminary Experiments}
\subsection{Experiment 1: Put Performance}
How many items you inserted, average latency or I/Os.

\subsection{Experiment 2: Get Performance}
Number of lookups, average read latency, any unoptimized vs. partial-optimized data.

\section{Challenges and Next Steps}
List major issues encountered (e.g., disk format, merges, tombstones). 
Mention at least three optimizations you plan to implement 
(e.g., dynamic Bloom filter bits-per-level, partial compaction, 
custom size ratio, etc.). 
Close with a short mention of concurrency plans for the final milestone.

\bibliographystyle{plain}
\begin{thebibliography}{1}
\bibitem{lsm1996}
  P. O'Neil, E. Cheng, etc. \textit{The Log-Structured Merge-Tree (LSM-tree)}. 1996.

\bibitem{leveldb}
  Google. \textit{LevelDB Documentation}, \url{https://github.com/google/leveldb}.

% add any references here
\end{thebibliography}

\end{document}
